\scanpage{65}
\mspage{64}

lieux (quand on n'est pas dans un atelier ensembliste tout au moins\ldots). Il n'est pas clair non plus, en général, que l'on ait, dans $\Sigma_{\unsure{\mathscr{F}|\mathcal{F}}}^\circ$,
\[
P\unsure{\wedge|\cap}\deleted{\unsure{\operatorname{Sup}_\Sigma|\operatorname{Sup}_{\mathscr{F}}}}\bigvee_i Q_i=\deleted{\unsure{\mathscr{F}|\Sigma|\operatorname{Sup}}}\bigvee_i(P\wedge Q_i)
\]
(distributivité de l'intersection des supports, par rapport aux ``sommes'' même finies\ldots).
Pourtant, cela a tendance à être vrai, de $\unsure{=|\overline m}$ que l'implication inverse de (x), dans les \deleted{\unsure{lieux où j'ai été}}\inserted{ateliers auxquels j'ai songé} jusqu'à présent, cf plus bas.

Tout \unsure{ceci|ce soit} se répéterait dans le cadre des \unsure{Stratos} $(M,\leq,\ll,\unsure{\gtrless|\lessgtr|\gtreqless|\lesseqgtr})$, \deleted{et} en termes de la seule relation $\parallel_M$ \deleted{et} de disjonction dans $M$. Et l'ens. $\Sigma_M^\circ$ des ``$M$-supports'' est en bij. canonique avec son homologue
\diagram{\begin{tikzcd}[ampersand replacement=\&]
\Sigma_M^\circ \arrow[r,shift left] \& \Sigma_{\mathscr{F}}^\circ \arrow[l,shift left]
\end{tikzcd}}
\margin{iso. d'ens. ord.}
(induit par les bijections de la p. 61)

% margin note written sideways; slanted upward
\margin{Si $A\subset\mathscr{F}$ est fermée pour $\leq$, on a\newline
$\operatorname{cosupp}_{\mathscr{F}}(A)\cap\unsure{L|M}=\operatorname{cosupp}_L(A\cap L)$,\newline
et pour $F\in\mathscr{F}$, on a\newline
$\operatorname{cosupp}_{\mathscr{F}}(F)\cap L=\operatorname{cosupp}_L\unsure{C^\circ(F)|\mathcal{C}^\circ(F)|(\mathring{\widetilde F})}$.}

commutant à l'opération ``complémentaire''.
Cette fois, le plus petit élément de $\Sigma_M^\circ$ est bien la partie \uline{vide} de $M$, et son plus grand élément est $M$.
